\chapter{Labeling and Target Construction}
\label{ch:labeling}

\begin{projectconnection}
In Phase~1 you build the label construction module (Task~6 in the implementation plan).
The choice of label is as important as the choice of features---the label \emph{is} your hypothesis.
Triple-barrier labeling with volatility-scaled thresholds is the primary labeling method;
meta-labeling adds a precision filter.
Your prediction targets (VIX innovations, realized vol, drawdowns) come from the
theoretical foundations in Chapters~\ref{ch:asset-pricing}--\ref{ch:intermediary}.
\end{projectconnection}

%% ═══════════════════════════════════════════════════════════════
\section{Why Labeling Matters}
\label{sec:labeling-why}

\begin{prereq}[Supervised Learning Recap]
Supervised learning maps inputs $\bX$ to labels $\by$.
The model minimizes $\loss(\hat{\by}, \by)$ over training data.
If $\by$ is poorly constructed---noisy, misaligned with the trading horizon, or
insensitive to the path of prices---no model can recover a useful signal.
Garbage in, garbage out applies with full force here.
\end{prereq}

Every supervised ML problem requires a label: the thing you predict.
In most ML applications the label is obvious (cat vs.\ dog, spam vs.\ not-spam).
In finance the label is a \emph{design choice}, and that choice encodes your hypothesis.

\textbf{The label defines the question.}
Predicting 5-day returns asks a different question than predicting whether
a stop-loss will be hit before a profit target.
The same features can yield a strong signal under one labeling scheme and
pure noise under another.

\subsection{Standard Return Labels}

The simplest label is the forward return over a fixed horizon:
\begin{equation}
\label{eq:fixed-horizon-label}
y_t \;=\; \frac{P_{t+h} - P_t}{P_t},
\end{equation}
where $h$ is a fixed holding period (1 day, 5 days, 21 days).
You then threshold: $y_t > \tau \Rightarrow +1$, $y_t < -\tau \Rightarrow -1$, else $0$.

Three problems:

\begin{enumerate}[leftmargin=2em]
\item \textbf{Path dependence is ignored.}
A position that hits $-5\%$ intraday and recovers to $+1\%$ at $t+h$
gets labeled $+1$, but no risk manager would have held through the drawdown.

\item \textbf{Fixed thresholds fail across asset classes.}
A $\pm 1\%$ threshold is reasonable for equities but trivial for
emerging-market FX and far too wide for G10 rates.
Heteroskedasticity within a single asset makes a fixed $\tau$ equally
inappropriate across time.

\item \textbf{Horizon selection is arbitrary.}
Why 5 days? Why not 3 or 10?
The choice of $h$ implicitly assumes a holding period that may not match
the signal's actual predictive horizon.
\end{enumerate}

\begin{intuition}[Labels as Strategy Blueprints]
Think of the label as a compressed description of what a trading strategy
would do.
A fixed-horizon return label describes a ``buy and hold for exactly $h$
days'' strategy---no stops, no targets, no early exit.
That is not how anyone trades.
The triple-barrier method (next section) labels data the way a
strategy actually behaves: entry, stop-loss, profit-target, and
maximum holding period.
\end{intuition}


%% ═══════════════════════════════════════════════════════════════
\section{Triple-Barrier Labeling}
\label{sec:triple-barrier}

The triple-barrier method, introduced by Lopez de Prado (2018, Ch.~3),
fixes every problem listed above.
It defines three barriers around each entry point and labels based on
which barrier the price path touches first.

\begin{definition}[Triple-Barrier Method]
Given an entry at time $t_0$ with price $P_{t_0}$:
\begin{enumerate}[leftmargin=2em]
\item \textbf{Upper barrier} (profit-taking):
$P_{t_0}\,(1 + m_u \cdot \sigma_t)$, where $m_u$ is the upper multiplier
and $\sigma_t$ is a rolling volatility estimate.
\item \textbf{Lower barrier} (stop-loss):
$P_{t_0}\,(1 - m_\ell \cdot \sigma_t)$, where $m_\ell$ is the lower multiplier.
\item \textbf{Vertical barrier} (expiration):
a maximum holding period $T_{\max}$ bars after $t_0$.
\end{enumerate}
The label is:
\[
y =
\begin{cases}
+1 & \text{if upper barrier is hit first,} \\
-1 & \text{if lower barrier is hit first,} \\
0  & \text{if vertical barrier is reached before either horizontal barrier.}
\end{cases}
\]
\end{definition}

\subsection{Volatility-Scaled Thresholds}

The barrier widths adapt to current market conditions through $\sigma_t$,
typically estimated as an exponentially weighted moving standard deviation
of log returns:
\begin{equation}
\label{eq:ewma-vol}
\sigma_t \;=\; \sqrt{\sum_{i=0}^{n-1} w_i \left(r_{t-i} - \bar{r}_t^{w}\right)^2},
\qquad w_i = \frac{(1-\lambda)\,\lambda^{i}}{1-\lambda^{n}},
\end{equation}
where $r_{t-i} = \ln P_{t-i} - \ln P_{t-i-1}$ are log returns,
$\bar{r}_t^{w} = \sum_i w_i\, r_{t-i}$ is the weighted mean,
and $\lambda = 2^{-1/\tau_h}$ for a half-life of $\tau_h$ days.
Lopez de Prado recommends a span of approximately 100 bars
(roughly $\tau_h \approx 50$ days for daily data).

\begin{prereq}[Exponentially Weighted Moving Average]
An EWMA gives more weight to recent observations.
The half-life $\tau_h$ is the number of periods after which a past
observation's weight drops to half.
A 20-day half-life means an observation from 20 days ago contributes
half as much as today's.
Pandas implements this as \texttt{df.ewm(halflife=20).std()}.
\end{prereq}

Why volatility scaling matters: a $2\%$ barrier is tight during a crisis
(daily vol $\approx 3\%$) but loose during calm markets (daily vol $\approx 0.5\%$).
Multiplying by $\sigma_t$ normalizes the barrier width in
``number of standard deviations,'' making labels comparable across time
and across asset classes.

\subsection{Triple-Barrier Diagram}

\begin{center}
\begin{tikzpicture}[xscale=1.1, yscale=1.4]
  % axes
  \draw[-{Stealth}] (0,0) -- (8.5,0) node[right] {\small time};
  \draw[-{Stealth}] (0,-2) -- (0,2.8) node[above] {\small price};

  % entry
  \filldraw (0.5,0) circle (2pt) node[below=3pt] {\footnotesize $P_{t_0}$};

  % upper barrier
  \draw[dashed, keyorange, thick] (0.5,1.6) -- (6.5,1.6)
    node[right, font=\footnotesize, keyorange] {upper: $P_{t_0}(1+m_u\sigma_t)$};

  % lower barrier
  \draw[dashed, warnred, thick] (0.5,-1.4) -- (6.5,-1.4)
    node[right, font=\footnotesize, warnred] {lower: $P_{t_0}(1-m_\ell\sigma_t)$};

  % vertical barrier
  \draw[dashed, defblue, thick] (6.5,-1.8) -- (6.5,2.2)
    node[above, font=\footnotesize, defblue] {$T_{\max}$};

  % price path (hits upper barrier)
  \draw[thick, black]
    (0.5,0) -- (1.2,0.3) -- (1.8,-0.2) -- (2.5,0.4) -- (3.0,0.1)
    -- (3.6,0.7) -- (4.1,0.5) -- (4.6,1.0) -- (5.1,1.2) -- (5.5,1.6);

  % hit marker
  \filldraw[keyorange] (5.5,1.6) circle (3pt);
  \node[above right, font=\footnotesize\bfseries, keyorange] at (5.5,1.7) {$y=+1$};

  % labels
  \draw[decorate, decoration={brace, amplitude=5pt, mirror}]
    (0.5,-2.0) -- (6.5,-2.0) node[midway, below=6pt, font=\footnotesize] {max holding period};
\end{tikzpicture}
\end{center}

\subsection{Worked Example: Step-by-Step Label Assignment}

\begin{workedexample}{Triple-Barrier Label Assignment}
\textbf{Setup.} Ten daily prices starting at $P_0 = 100$.
Rolling volatility (EWMA, 20-day half-life) at entry: $\sigma_0 = 1.0\%$.
Multipliers: $m_u = m_\ell = 2$.
Maximum holding period: $T_{\max} = 5$ days.

\textbf{Barriers.}
Upper: $100 \times (1 + 2 \times 0.01) = 102.00$.
Lower: $100 \times (1 - 2 \times 0.01) = 98.00$.
Vertical: day 5.

\textbf{Price path.}
\begin{center}
\begin{tabular}{lccccc}
\toprule
Day & 1 & 2 & 3 & 4 & 5 \\
\midrule
Price & 100.40 & 100.80 & 101.50 & 101.20 & 102.10 \\
Hit? & --- & --- & --- & --- & upper \\
\bottomrule
\end{tabular}
\end{center}

On day 5 the price crosses 102.00 for the first time $\Rightarrow$
upper barrier is hit before the vertical barrier $\Rightarrow$ $y = +1$.

\textbf{Alternate path.}
If instead the prices were $\{99.80, 99.30, 98.50, 97.90, \ldots\}$,
the price crosses 98.00 on day 3 $\Rightarrow$ $y = -1$.

\textbf{Alternate path 2.}
Prices $\{100.20, 99.80, 100.10, 100.50, 100.30\}$.
Neither barrier is breached within 5 days.
At the vertical barrier: $P_5 = 100.30 > P_0 = 100$, so $y = 0$
(or optionally $y = \operatorname{sign}(P_5 - P_0) = +1$, depending
on configuration).
\end{workedexample}

\begin{keyidea}[Why Triple-Barrier Labeling Works]
Triple-barrier labeling aligns the label with how a trading strategy
would actually behave: a position has a stop-loss, a profit target,
and a maximum holding period.
The label records \emph{which constraint binds first}, not just the
terminal return.
This produces labels that are more informative and less noisy than
fixed-horizon returns, because the label reflects the path of prices,
not just the endpoint.
\end{keyidea}


%% ═══════════════════════════════════════════════════════════════
\section{Meta-Labeling}
\label{sec:meta-labeling}

Triple-barrier labeling tells you the \emph{outcome} of a trade.
Meta-labeling (Lopez de Prado 2018, Ch.~3.6) adds a second layer:
a model that predicts \emph{whether to take the trade at all}.

\begin{definition}[Meta-Labeling]
A two-stage labeling procedure:
\begin{enumerate}[leftmargin=2em]
\item \textbf{Primary model} predicts the \emph{side}: long ($+1$) or short ($-1$).
This can be any model---ML, fundamental, or rule-based
(e.g., moving-average crossover, carry signal, mean-reversion filter).
\item \textbf{Secondary (meta) model} predicts whether the primary model is
\emph{correct}: $1$ if the trade is profitable (per triple-barrier), $0$ otherwise.
\end{enumerate}
The meta-label is:
\[
y^{\text{meta}}_t =
\begin{cases}
1 & \text{if primary side} \times \text{triple-barrier label} > 0, \\
0 & \text{otherwise.}
\end{cases}
\]
\end{definition}

\subsection{Precision--Recall Tradeoff}

The primary model is designed for \textbf{high recall}: it flags most
profitable opportunities, even at the cost of many false positives.
The meta-model adds \textbf{precision}: it filters out trades where
the primary model is likely wrong.

\begin{prereq}[Precision and Recall]
\textbf{Precision} $= \text{TP} / (\text{TP} + \text{FP})$: of all trades the model takes,
what fraction are profitable?
\textbf{Recall} $= \text{TP} / (\text{TP} + \text{FN})$: of all profitable opportunities,
what fraction does the model capture?
A model that trades on every signal has recall $= 1$ but may have
low precision.
Meta-labeling sacrifices some recall for substantially higher precision.
\end{prereq}

In empirical tests using synthetic data and equity strategies, adding
meta-labeling to a Bollinger Band mean-reversion strategy increased
out-of-sample precision from 0.17 to 0.20 and accuracy from 17\% to
63\% (Singh and Joubert, Hudson \& Thames research note).
For a trend-following SMA crossover, out-of-sample precision rose from
0.48 to 0.54.
These gains come from filtering false positives, not from discovering
new signal.

\subsection{Betting Size from Meta-Probabilities}

The meta-model outputs a probability $p_t \in [0,1]$.
This directly maps to position size:
\begin{equation}
\label{eq:bet-size}
\text{size}_t \;=\; f(p_t) \;\times\; \text{side}_t,
\end{equation}
where $f$ is a monotone mapping (e.g., $f(p) = 2p - 1$ so that
$p = 0.5 \Rightarrow$ no bet, $p = 1.0 \Rightarrow$ full position).
The side comes from the primary model; the magnitude comes from the
meta-model's confidence.

\begin{warning}[Meta-Labeling Requires a Reasonable Primary Model]
Meta-labeling filters false positives from the primary model.
If the primary model is random (accuracy $\approx 50\%$ with no structure),
there is nothing systematic for the meta-model to learn.
Always verify that the primary model has above-chance directional accuracy
before investing effort in meta-labeling.
A ridge regression baseline (Chapter~\ref{ch:regularized}) is a
reasonable first primary model.
\end{warning}


%% ═══════════════════════════════════════════════════════════════
\section{Label Uniqueness and Concurrency}
\label{sec:label-uniqueness}

\begin{prereq}[Why i.i.d.\ Assumptions Fail for Overlapping Labels]
Standard ML assumes training samples are independent and identically
distributed (i.i.d.).
Triple-barrier labels with a 5-day maximum holding period create
overlapping return windows: the label for day $t$ depends on prices
$\{P_{t+1}, \ldots, P_{t+5}\}$, and the label for day $t+1$ depends on
$\{P_{t+2}, \ldots, P_{t+6}\}$.
Four of five days overlap.
This makes consecutive labels highly correlated, violating the
i.i.d.\ assumption.
Ignoring this inflates apparent model performance, because the model
``memorizes'' shared return information rather than learning genuine patterns.
\end{prereq}

\subsection{Concurrency and Uniqueness}

Lopez de Prado (2018, Ch.~4) quantifies this overlap with two
measures.

\begin{definition}[Concurrency Count]
For a set of $N$ labeled events, each spanning a time window
$[t_i^{\text{start}}, t_i^{\text{end}}]$, the concurrency at bar $t$ is:
\begin{equation}
\label{eq:concurrency}
c_t \;=\; \sum_{i=1}^{N} \mathbf{1}\!\left[t_i^{\text{start}} \leq t \leq t_i^{\text{end}}\right],
\end{equation}
where $\mathbf{1}[\cdot]$ is the indicator function.
$c_t$ counts how many labels are ``active'' at time $t$---i.e., how
many overlapping return windows include bar $t$.
\end{definition}

\begin{definition}[Average Uniqueness]
The average uniqueness of label $i$ is:
\begin{equation}
\label{eq:avg-uniqueness}
\bar{u}_i \;=\; \frac{1}{|T_i|}\sum_{t \in T_i} \frac{1}{c_t},
\end{equation}
where $T_i = \{t : t_i^{\text{start}} \leq t \leq t_i^{\text{end}}\}$
is the set of bars spanned by label $i$.
A label with $\bar{u}_i = 1$ has no overlap with any other label.
A label with $\bar{u}_i = 0.25$ shares its information with roughly
three other concurrent labels on average.
\end{definition}

\begin{workedexample}{Computing Average Uniqueness}
Three events with the following spans:
\begin{center}
\begin{tabular}{lccc}
\toprule
Event & Start & End & Bars \\
\midrule
$A$ & 1 & 3 & $\{1,2,3\}$ \\
$B$ & 2 & 4 & $\{2,3,4\}$ \\
$C$ & 4 & 5 & $\{4,5\}$ \\
\bottomrule
\end{tabular}
\end{center}

Concurrency at each bar:
$c_1 = 1$, $c_2 = 2$, $c_3 = 2$, $c_4 = 2$, $c_5 = 1$.

Average uniqueness:
\begin{align*}
\bar{u}_A &= \tfrac{1}{3}\left(\tfrac{1}{1} + \tfrac{1}{2} + \tfrac{1}{2}\right) = \tfrac{1}{3}(2.0) \approx 0.667, \\
\bar{u}_B &= \tfrac{1}{3}\left(\tfrac{1}{2} + \tfrac{1}{2} + \tfrac{1}{2}\right) = 0.500, \\
\bar{u}_C &= \tfrac{1}{2}\left(\tfrac{1}{2} + \tfrac{1}{1}\right) = 0.750.
\end{align*}
Event $C$ is most unique; event $B$ is least unique.
In sample-weighted training, $C$ receives the highest weight.
\end{workedexample}

\subsection{Using Uniqueness for Sample Weighting}

There are two practical applications:

\begin{enumerate}[leftmargin=2em]
\item \textbf{Sample weights.}
Set $w_i = \bar{u}_i$ and pass these weights to the model's loss
function.
Both LightGBM and scikit-learn accept a \texttt{sample\_weight}
argument.
High-uniqueness samples contribute more to the gradient;
low-uniqueness (highly overlapping) samples are downweighted.

\item \textbf{Bagging adjustment.}
When using bagged models (random forests, bagged trees), set
\texttt{max\_samples}$\;= \overline{u}$, the average uniqueness across
all labels.
This prevents bootstrap samples from containing many overlapping
observations.
Lopez de Prado further proposes \emph{sequential bootstrapping},
which draws samples one at a time and reweights the sampling
distribution after each draw to penalize overlap with
already-selected observations.
\end{enumerate}


%% ═══════════════════════════════════════════════════════════════
\section{Prediction Targets for the Project}
\label{sec:prediction-targets}

The labeling method (triple-barrier, meta-labeling) defines \emph{how}
you label.
The prediction target defines \emph{what} you label.
This section specifies the four targets from the design spec and
explains why each is motivated by theory.

\subsection{VIX Innovations}

\begin{definition}[VIX Innovation]
\begin{equation}
\label{eq:vix-innovation}
\Delta \text{VIX}_t \;=\; \text{VIX}_t \;-\; \text{VIX}_{t-1}.
\end{equation}
This is the daily change in implied volatility.
Positive $\Delta\text{VIX}_t$ means the market's expectation of
future volatility increased.
\end{definition}

\begin{keyresult}[Adrian and Shin (2010): Dealer Repos Forecast VIX]
Adrian and Shin show that changes in U.S.\ primary dealer repo positions
forecast innovations in the VIX index.
A forecasting regression of VIX changes on lagged repo growth produces
$R^2 = 11.6\%$, compared to $R^2 = 8.9\%$ using only the lagged VIX
level---a relative improvement of 30\%.
The predictive power comes primarily through the \emph{volatility risk
premium} (the gap between implied and realized volatility), not through
actual volatility changes.
Interpretation: when dealers expand their balance sheets
(repo growth $> 0$), risk appetite rises and implied volatility
compresses; when dealers contract (repo growth $< 0$), risk appetite
falls and implied volatility spikes.

\medskip
\noindent
\textit{Source:} Adrian, T.\ and Shin, H.S.\ (2010),
``Liquidity and Leverage,''
\textit{Journal of Financial Intermediation}, 19(3), 418--437.
\end{keyresult}

Why this is the cleanest target: VIX innovations are a single,
tradeable, observable time series with a direct theoretical link to
the risk data you have in SecDB.
Adrian--Shin used stale quarterly Fed data; SecDB gives you \emph{daily}
dealer-level risk outputs.
If the quarterly signal works, the daily signal should be strictly
stronger.

\subsection{Realized Volatility}

\begin{definition}[Realized Volatility]
\begin{equation}
\label{eq:realized-vol}
\text{RV}_t \;=\; \sqrt{\sum_{i=1}^{n} r_{t,i}^2},
\end{equation}
where $r_{t,i}$ is the $i$-th intraday return on day $t$ and $n$ is
the number of intraday intervals.
For daily data without intraday granularity, the simpler proxy is
$\text{RV}_t = |r_t|$ or $\text{RV}_t = r_t^2$.
\end{definition}

\begin{prereq}[Why Volatility Is Easier to Predict Than Returns]
Returns have very low signal-to-noise ratios (SNR).
A Sharpe ratio of 1.0 on daily returns corresponds to SNR $\approx
1/\sqrt{252} \approx 0.063$---the signal is 6\% of the noise.
Volatility is more persistent (high autocorrelation), more
predictable (GARCH models capture 30--50\% of variance), and
does not require predicting direction.
Predicting $\text{RV}_{t+1}$ is a strictly easier problem than
predicting $r_{t+1}$.
\end{prereq}

Realized volatility is the second-priority target.
It is useful both as a standalone prediction (vol-timing strategies)
and as an input to volatility-adjusted position sizing.
Predicting it with risk-cube features tests whether SecDB data
captures information about future realized risk beyond what
standard GARCH models provide.

\subsection{Drawdowns by Most-Concentrated Asset Class}

\begin{definition}[Maximum Drawdown over Horizon]
\begin{equation}
\label{eq:drawdown}
\text{DD}_{t,h} \;=\; \max_{0 \leq s \leq h}\;
\frac{\max_{0 \leq u \leq s} P_{t+u} - P_{t+s}}{\max_{0 \leq u \leq s} P_{t+u}},
\end{equation}
where $h$ is the forward-looking horizon.
This measures the worst peak-to-trough decline in the window $[t, t+h]$.
\end{definition}

The design spec identifies factor concentration ($\HHI$ of VaR
decomposition; see Chapter~\ref{ch:features}) as a priority feature.
The natural target to pair it with is drawdowns in the most-concentrated
asset class: when risk is concentrated and then unwinds (fire sales
per Coval and Stafford 2007, as discussed in Chapter~\ref{ch:intermediary}),
the concentrated asset class should experience the largest drawdown.

This target is noisier than VIX innovations.
Use it as a secondary target, not a starting point.

\subsection{Cross-Asset Momentum Reversals}

\begin{definition}[Momentum Reversal Indicator]
Label $y_t = +1$ if the trailing 12-month winner underperforms
the trailing 12-month loser over the next $h$ days (reversal), and
$y_t = -1$ otherwise (continuation).
Compute across asset classes (equities, rates, credit, FX)
rather than within.
\end{definition}

Cross-asset momentum reversals are the most complex target.
The hypothesis: when VaR utilization is near the limit, forced
deleveraging causes recent winners to reverse as dealers unwind
their most profitable (and most concentrated) positions.
He, Kelly, and Manela (2017, discussed in Chapter~\ref{ch:intermediary})
show that a single intermediary SDF prices assets across classes,
so cross-asset reversals should be predictable from cross-asset
risk data.

This target has the longest horizon and the fewest independent
observations.
Reserve it for Phase~4A if VIX innovations and realized vol
produce positive results.

\begin{keyidea}[Target Priority Order]
Start with VIX innovations (cleanest target per Adrian--Shin, single
observable series, direct theory link) and realized volatility
(higher SNR than returns, more predictable).
Only move to drawdowns and momentum reversals after establishing
that SecDB features contain signal on the easier targets.
If VIX innovations and realized vol both show zero IC after purged
cross-validation (Chapter~\ref{ch:validation}), the harder targets
will not rescue the project---that is the Week~13 checkpoint.
\end{keyidea}


%% ═══════════════════════════════════════════════════════════════
\section{Practical Considerations}
\label{sec:labeling-practical}

\subsection{Choosing Barrier Parameters}

The triple-barrier method has three parameters: $m_u$, $m_\ell$, and $T_{\max}$.
Guidelines from practice:

\begin{itemize}[leftmargin=2em]
\item \textbf{Symmetric barriers} ($m_u = m_\ell$) are the default.
Start with $m_u = m_\ell = 2$ (i.e., barriers at $\pm 2\sigma$).
Asymmetric barriers encode a view that profits and losses are
treated differently---only use this if you have a specific reason.

\item \textbf{Maximum holding period.}
Match $T_{\max}$ to the signal's expected decay.
For daily risk-cube features, $T_{\max} = 5$ to $21$ trading days is
reasonable.
Too short and most labels hit the vertical barrier ($y=0$);
too long and labels become noisy.

\item \textbf{Sensitivity analysis.}
Always test at least 3 parameter combinations
(e.g., $m = \{1.5, 2.0, 2.5\}$, $T_{\max} = \{5, 10, 21\}$)
and report results for all.
A signal that only appears at one parameter setting is fragile.
Log every combination in the experiment tracker for honest Deflated
Sharpe Ratio accounting (Chapter~\ref{ch:validation}).
\end{itemize}

\subsection{Handling Label Imbalance}

Triple-barrier labels often produce imbalanced classes.
In low-volatility periods, most labels hit the vertical barrier
($y = 0$), creating a dominant class.
Options:

\begin{enumerate}[leftmargin=2em]
\item \textbf{Merge} $y=0$ into the sign of the terminal return,
reducing to a binary $\{-1, +1\}$ problem.
\item \textbf{Upsample} the minority class with replacement before training.
\item \textbf{Class weights.} Set
$w_c = N / (k \cdot N_c)$, where $N$ is total samples, $k$ is number
of classes, and $N_c$ is class count.
Scikit-learn and LightGBM support this natively.
\end{enumerate}


%% ═══════════════════════════════════════════════════════════════
\section{Summary}
\label{sec:labeling-summary}

\begin{center}
\begin{tabular}{p{3.8cm} p{9.5cm}}
\toprule
\textbf{Concept} & \textbf{Key Takeaway} \\
\midrule
Fixed-horizon labels &
Simple but ignore path, fail across asset classes, and assume
an arbitrary holding period. \\[4pt]
Triple-barrier labeling &
Three barriers (profit, stop, time) with volatility-scaled thresholds.
Label $= $ which barrier is hit first.
Aligns labels with actual strategy behavior. \\[4pt]
Meta-labeling &
Primary model predicts side; secondary model predicts whether primary
is correct. Trades recall for precision; meta-probability $\to$
position size. \\[4pt]
Label uniqueness &
Overlapping return windows violate i.i.d.
Average uniqueness $\bar{u}_i = \frac{1}{|T_i|}\sum_{t \in T_i} 1/c_t$
gives sample weights for training. \\[4pt]
VIX innovations &
Cleanest target. Adrian--Shin (2010): dealer repo growth predicts
$\Delta$VIX with $R^2$ improvement from 8.9\% to 11.6\%. \\[4pt]
Realized volatility &
Higher SNR than returns. Tests whether SecDB data beats GARCH. \\[4pt]
Drawdowns &
Pair with factor concentration ($\HHI$). Noisier; use as secondary target. \\[4pt]
Momentum reversals &
Cross-asset, longest horizon. Reserve for Phase 4A. \\
\bottomrule
\end{tabular}
\end{center}

The labeling choices made here flow directly into the validation
framework (Chapter~\ref{ch:validation}): purged cross-validation must
respect the overlapping windows created by triple-barrier labels,
and the Deflated Sharpe Ratio must account for every target--parameter
combination tested.
Build the label module (Task~6) before any signal testing, and
commit to the experiment log from the first run.
